\documentclass[11pt,a4paper]{article}
\usepackage[italian]{babel}
\usepackage[margin=2.5cm]{geometry}
\usepackage{parskip}
\usepackage{amsmath, amssymb, amsthm}
\usepackage{xcolor}
\usepackage{listings}
\usepackage{tikz}
\usepackage[most]{tcolorbox}
\usepackage{hyperref}

% --- Riquadri con bordo nero, senza numero ---
\tcbset{nota/.style={colback=white, colframe=black, boxrule=0.5pt,
  sharp corners}}
\NewTColorBox{definizione}{}{nota,
  before upper={\textbf{Definizione.}\ }}
\NewTColorBox{osservazione}{}{nota,
  before upper={\textbf{Osservazione.}\ }}
\theoremstyle{definition}
\newtheorem*{esempio}{Esempio}
\newtheorem*{esercizio}{Esercizio}

% --- Codice C ---
% Uso: \begin{codice} ... \end{codice}; \begin{risultato} ... \end{risultato}
%      \lstinline|printf("ciao");| per il codice dentro al testo
\lstdefinestyle{c}{
  language=C,
  basicstyle=\ttfamily\small,
  keywordstyle=\color{blue}\bfseries,
  commentstyle=\color{gray}\itshape,
  stringstyle=\color{red!70!black},
  numbers=left,
  numberstyle=\tiny\color{gray},
  numbersep=8pt,
  columns=flexible,
  keepspaces=true,
  breaklines=true,
  showstringspaces=false,
  tabsize=4,
  morekeywords={include, define},
  literate={à}{{\`a}}1 {è}{{\`e}}1 {é}{{\'e}}1 {ì}{{\`i}}1
           {ò}{{\`o}}1 {ù}{{\`u}}1
}
\lstset{style=c}
\newtcblisting{codice}{listing only, nota, left=6mm,
  listing options={style=c}}
\newtcblisting{risultato}{listing only, colback=black!5, colframe=black,
  boxrule=0.5pt, sharp corners,
  listing options={basicstyle=\ttfamily\small, numbers=none, language={}}}

% --- Diagrammi di flusso (simboli standard) ---
\usetikzlibrary{shapes.geometric, arrows.meta, positioning}
\tikzset{
  terminale/.style = {ellipse, draw, minimum width=2.5cm, minimum height=0.9cm},
  io/.style        = {trapezium, trapezium left angle=70, trapezium right angle=110,
                      draw, minimum width=2.5cm, minimum height=0.9cm},
  processo/.style  = {rectangle, draw, minimum width=2.5cm, minimum height=0.9cm},
  decisione/.style = {diamond, draw, aspect=2, minimum width=2.5cm},
  freccia/.style   = {thick, -{Stealth}}
}

\title{Errori di compilazione e di esecuzione}
\author{Marco Cerbella}
\date{Lezione 5 -- 30 settembre 2026}

\begin{document}
\maketitle

\section{Compile-time e run-time}

La vita di un programmatore ha due momenti: il \emph{compile-time}, quando si
esegue \texttt{gcc}, e il \emph{run-time}, quando si esegue il programma.

\begin{definizione}
Un \emph{errore di compilazione} (\emph{compile-time}) è un errore rilevato dal
compilatore. Cause comuni: errori di sintassi, come un punto e virgola mancante
o l'uso di una parola riservata (ad esempio \lstinline|while|) come nome.
\end{definizione}

\begin{definizione}
Un \emph{errore di esecuzione} (\emph{run-time}) è un errore del programma che
si verifica mentre il programma è in esecuzione.
\end{definizione}

\section{Esempio di errore di compilazione}

Nel programma del MCD manca il \texttt{;} dopo la seconda \lstinline|scanf|:
\begin{codice}
printf("Enter first positive integer: \n");
scanf("%d", &b)        // manca il ';'
\end{codice}

Il compilatore segnala:
\begin{risultato}
$ gcc -o euclid euclid.c
euclid.c:9:20: error: expected ';' after expression
   scanf("%d", &b)
                 ^
                 ;
1 error generated.
\end{risultato}

Il messaggio indica file, riga e colonna dell'errore: bisogna imparare a
leggere l'output del compilatore per eliminare tutti gli errori di compilazione.

\section{Esempio di errore di esecuzione}

Il programma compila senza problemi, ma dereferenzia un puntatore nullo:
\begin{codice}
#include <stdio.h>
int main() {
    int *a = NULL;

    printf("%d", *a);    // accesso a memoria non valida

    int i = 0;
    while (i < 4) {
        printf("Hello World");
        i++;
    }
    return 0;
}
\end{codice}
\begin{risultato}
$ gcc -o runtime_error runtime_error.c
$ ./runtime_error
Segmentation fault: 11
\end{risultato}

\begin{osservazione}
Eliminare gli errori di compilazione non basta: bisogna testare il programma
con molti input diversi per eliminare (il più possibile) gli errori di
esecuzione.
\end{osservazione}

\begin{center}
\emph{``Testing shows the presence, not the absence of bugs''} -- E.\,W.\ Dijkstra
\end{center}
\end{document}
